\section{Proofs for Section~\ref{sec:labor}}\label{app:proofs-labor}

Notation is as in Section~\ref{sec:labor}, and $(x)_+=\max\{x,0\}$. Times are measured from date 0 of Assumption~\ref{ass:labor-proc} unless stated otherwise.

\paragraph{Proposition~\ref{prop:labor-lead}.}
(a) By Assumption~\ref{ass:labor-proc}(iii) the parallel run in $\mathfrak c$ occupies $[T^{\rm st},t^{\rm gate})$ with $t^{\rm gate}=T^{\rm st}+T_{\rm par}$, and the legacy mode cannot be retired before the run ends. By part (i), dismissals covered by a procedure that starts at $t^A$ cannot take effect before $t^A+T^{\rm leg}$. The earliest cutover is the later of the two dates, $t^{\rm dis}=\max\{t^{\rm gate},t^A+T^{\rm leg}\}$. Without legal constraints ($T^{\rm st}=0$ and $T^{\rm leg}=0$) the saving would start at $T_{\rm par}$. With $t^A=0$, $t^{\rm dis}-T_{\rm par}=\max\{t^{\rm gate}-T_{\rm par},T^{\rm leg}-T_{\rm par}\}=\max\{T^{\rm st},T^{\rm leg}-T_{\rm par}\}$. With $t^A=t^{\rm gate}$, $t^{\rm dis}=t^{\rm gate}+T^{\rm leg}$, so $t^{\rm dis}-T_{\rm par}=T^{\rm st}+T^{\rm leg}$.

(b) If all jurisdictions cut over together at $\bar t:=\max_{\mathfrak c}t^{\rm dis}_{\mathfrak c}$, jurisdiction $\mathfrak c$ waits $\bar t-t^{\rm dis}_{\mathfrak c}$ beyond its own earliest date, which to first order forgoes $W_{\mathfrak c}(\bar t-t^{\rm dis}_{\mathfrak c})$. A division into streams $\mathfrak C_1,\dots,\mathfrak C_q$, each cutting over at $T_i:=\max_{\mathfrak c\in\mathfrak C_i}t^{\rm dis}_{\mathfrak c}$, costs $\sum_i\sum_{\mathfrak c\in\mathfrak C_i}W_{\mathfrak c}(T_i-t^{\rm dis}_{\mathfrak c})+(q-1)O_w$. Take an optimal division, and suppose some $\mathfrak c\in\mathfrak C_i$ has $t^{\rm dis}_{\mathfrak c}\le T_j<T_i$ for some other stream $\mathfrak C_j$. Moving $\mathfrak c$ to $\mathfrak C_j$ leaves $T_j$ unchanged, does not raise $T_i$, lowers $\mathfrak c$'s wait from $T_i-t^{\rm dis}_{\mathfrak c}$ to $T_j-t^{\rm dis}_{\mathfrak c}$, and does not change the number of streams unless $\mathfrak C_i$ empties, which saves $O_w$. The cost does not rise. Repeating the move gives an optimal division in which every member of a stream has $t^{\rm dis}$ above the cutover date of every earlier stream, so the streams are intervals in the order of $t^{\rm dis}$. For two streams, with the fastest group $\mathfrak C_1$ cutting over at $T_1=\max_{\mathfrak C_1}t^{\rm dis}_{\mathfrak c}$ and the rest at $\bar t$, the saving relative to full synchronization is $\sum_{\mathfrak C_1}W_{\mathfrak c}(\bar t-T_1)$ and the extra cost is $O_w$, which gives the condition.

(c) Let the dismissal dates be $t_1\le\dots\le t_{N^{\rm dis}}$ and $n^{\rm b}:=N^{\rm thr}-1$. Staying below the threshold in every window means no $n^{\rm b}+1$ consecutive dismissals fall within a window, so $t_{i+n^{\rm b}}-t_i\ge T^{\rm win}$ for every $i\le N^{\rm dis}-n^{\rm b}$. With $i^{\rm b}:=\lceil N^{\rm dis}/n^{\rm b}\rceil-1$ we have $1+i^{\rm b}n^{\rm b}\le N^{\rm dis}$, and summing over $i=1,1+n^{\rm b},\dots,1+(i^{\rm b}-1)n^{\rm b}$ gives $t_{1+i^{\rm b}n^{\rm b}}-t_1\ge i^{\rm b}T^{\rm win}$, so $t_{N^{\rm dis}}-t_1\ge i^{\rm b}T^{\rm win}$. Placing blocks of $n^{\rm b}$ dismissals at $0,T^{\rm win},2T^{\rm win},\dots$ attains the bound.

(d) Write $T:=T^{\rm leg}$ and $\tilde t:=t^{\rm gate}$. If $t^A=\tilde t$, the procedure starts only after a success, so the expected cost is $C(\tilde t)=\varpi(c^{\rm ann}+W)T$. For $t^A<\tilde t$ the procedure is announced before the outcome is known: on success, the announced state lasts $t^{\rm dis}-t^A$ and the saving waits $t^{\rm dis}-\tilde t$; on failure, the procedure is withdrawn at the gate. Hence
\[
C(t^A)=\varpi\bigl[c^{\rm ann}(t^{\rm dis}-t^A)+W(t^{\rm dis}-\tilde t)\bigr]+(1-\varpi)\bigl[c^{\rm ann}(\tilde t-t^A)+C^{\rm wd}\bigr],\qquad t^{\rm dis}=\max\{\tilde t,t^A+T\}.
\]
On $[0,(\tilde t-T)_+]$, $t^{\rm dis}=\tilde t$ and $C=c^{\rm ann}(\tilde t-t^A)+(1-\varpi)C^{\rm wd}$, which decreases in $t^A$. On $[(\tilde t-T)_+,\tilde t)$, $t^{\rm dis}=t^A+T$ and $C$ is affine in $t^A$ with slope $\varpi W-(1-\varpi)c^{\rm ann}$, and its limit as $t^A\to\tilde t$ is $\varpi(c^{\rm ann}+W)T+(1-\varpi)C^{\rm wd}\ge C(\tilde t)$. An affine function attains its minimum over an interval at an endpoint, so the minimum over $[0,\tilde t]$ is attained at $(\tilde t-T)_+$ or at $\tilde t$. If $T\le\tilde t$, $C((\tilde t-T)_+)=c^{\rm ann}T+(1-\varpi)C^{\rm wd}$ and $C(\tilde t)-C(\tilde t-T)=\varpi WT-(1-\varpi)(c^{\rm ann}T+C^{\rm wd})$. If $T>\tilde t$, $C(0)=\varpi[c^{\rm ann}T+W(T-\tilde t)]+(1-\varpi)(c^{\rm ann}\tilde t+C^{\rm wd})$ and $C(\tilde t)-C(0)=\varpi W\tilde t-(1-\varpi)(c^{\rm ann}\tilde t+C^{\rm wd})$. In both cases $C(\tilde t)<C((\tilde t-T)_+)$ if and only if \eqref{eq:labor-ann} holds with $T^{\rm ov}=\min\{T,\tilde t\}$, and solving for $\varpi$ gives $\varpi^A$. Under the bound $t^A\le t^{\rm cont}<\tilde t$, $C$ is piecewise affine on $[0,t^{\rm cont}]$ with a single possible kink at $\tilde t-T$, so its minimum is at $0$, at $t^{\rm cont}$, or at $\tilde t-T$ if that lies in between.

(e) Measure time from the start of the run. Ending the run at $T$ means the firm detects a failure during it with probability $1-e^{-T/T_{\rm det}}$ (Assumption~\ref{ass:tasks-transition}(iv)). Before $T^{\rm hold}$ the legacy staff cannot be dismissed, so the duplication premium $\omega_{\rm par}C_0$ is paid whether or not the run continues, and no saving is available, so none is forgone; continuing the run costs only $\hat c'$ a year. After $T^{\rm hold}$ a year of running costs $\hat c$, as in Proposition~\ref{prop:tasks-transition}(d). The part of the undiscounted NPV that depends on $T$ is therefore
\[
J(T)=-\hat c'\min\{T,T^{\rm hold}\}-\hat c\,(T-T^{\rm hold})_+-(1-\varpi_p)e^{-T/T_{\rm det}}(C^{\rm str}+L_{\rm roll}).
\]
Since $\hat c'\le\hat c$, the first two terms form a concave function, and the last is concave because $e^{-T/T_{\rm det}}$ is convex; so $J$ is concave. Its right derivative is $-\hat c'+\mathfrak b(T)$ for $T<T^{\rm hold}$ and $-\hat c+\mathfrak b(T)$ for $T\ge T^{\rm hold}$, where $\mathfrak b(T):=(1-\varpi_p)(C^{\rm str}+L_{\rm roll})e^{-T/T_{\rm det}}/T_{\rm det}$ is decreasing. $T^{**}_{\rm par}$ and $T^*_{\rm par}$ are the points where $\mathfrak b$ equals $\hat c'$ and $\hat c$ (or zero if $\mathfrak b(0)$ is below them), so $T^*_{\rm par}\le T^{**}_{\rm par}$. If $T^{**}_{\rm par}\le T^{\rm hold}$, $J$ peaks at $T^{**}_{\rm par}$, which equals the right-hand side of \eqref{eq:labor-freerun} because $\max\{T^{\rm hold},T^*_{\rm par}\}\ge T^{\rm hold}\ge T^{**}_{\rm par}$. If $T^{**}_{\rm par}>T^{\rm hold}$, $J$ increases on $[0,T^{\rm hold}]$, and on $[T^{\rm hold},\infty)$ it peaks at $\max\{T^{\rm hold},T^*_{\rm par}\}\le T^{**}_{\rm par}$, which again equals the right-hand side. $\square$

\paragraph{Proposition~\ref{prop:labor-survivor}.}
(a) A survivor present at the announcement quits within $T_s$ with probability $1-e^{-(\nu^q+\nu^s)T_s}$, against $1-e^{-\nu^qT_s}$ without it. The difference is $e^{-\nu^qT_s}(1-e^{-\nu^sT_s})\le\nu^sT_s$, since $1-e^{-x}\le x$. Summing over the $N^{\rm sv}$ survivors gives $N^{\rm q}$. Each extra quit costs $c^{\rm rep}w$, and output falls by $\delta^{\rm eff}wN^{\rm sv}$ a year for $T_s$ years, which gives $C^{\rm surv}$.

(b) Let $\nu^s(x)/x$ be nonincreasing on $(0,1]$. For $x,y>0$ with $x+y\le1$, $\nu^s(x+y)=x\,\nu^s(x+y)/(x+y)+y\,\nu^s(x+y)/(x+y)\le x\,\nu^s(x)/x+y\,\nu^s(y)/y=\nu^s(x)+\nu^s(y)$, and by induction $\nu^s(\sum_i\mathsf d_i)\le\sum_i\nu^s(\mathsf d_i)$. To first order (with the survivor base held at $N^{\rm sv}$ across tranches), the extra quits are $N^{\rm sv}T_s\sum_i\nu^s(\mathsf d_i)\ge N^{\rm sv}T_s\nu^s(\mathsf d)$. For equal tranches, $n\nu^s(\mathsf d/n)=\mathsf d\cdot\nu^s(\mathsf d/n)/(\mathsf d/n)$, and since $\mathsf d/n$ decreases in $n$ and $\nu^s(x)/x$ is nonincreasing, this is nondecreasing in $n$. With $\nu^s(x)=\eta_0\mathbf1\{x>0\}+\eta_1x$ and $\delta^{\rm eff}$ of the same form, the first-order cost of the program is $wN^{\rm sv}T_s\sum_i[c^{\rm rep}(\eta_0+\eta_1\mathsf d_i)+\delta_0+\delta_1\mathsf d_i]$ over the announced waves $i$. Since $\sum_i\mathsf d_i$ is fixed, the linear terms do not depend on the number of waves $n_w$, and the rest equals $n_wO_s$. So the survivor cost adds $n_wO_s$ to \eqref{eq:ext-progcost}, exactly as the overhead $O_w$ enters it. With equal waves of $n=N/n_w$ processes the cost of Proposition~\ref{prop:ext-waves}(b) becomes $T_{\rm par}[WN(n_w+1)/2+c^{\rm dr}_1N+c^{\rm dr}_2N^2/n_w]+n_w(O_w+O_s)$, whose first-order condition in $n_w$ gives the wave size $N/n_w=\sqrt{[WN/2+(O_w+O_s)/T_{\rm par}]/c^{\rm dr}_2}$.

(c) By Assumption~\ref{ass:labor-survivor}(iii), each holder's excess departures arrive at the rate $(1-\kappa_{\rm ret})\nu^x$ during the announced part of the run, independently of the baseline reflected in $\varpi^0_p$ and of the other holders. The holder causes no failure if no excess departure occurs within $\tau^{\rm ann}$, probability $e^{-(1-\kappa_{\rm ret})\nu^x\tau^{\rm ann}}$, or if a departure occurs but its context is captured, probability $(1-e^{-(1-\kappa_{\rm ret})\nu^x\tau^{\rm ann}})\kappa_{\rm doc}$. The sum is $\kappa_{\rm doc}+(1-\kappa_{\rm doc})e^{-(1-\kappa_{\rm ret})\nu^x\tau^{\rm ann}}$, and independence across the $n_K$ holders gives the equality in \eqref{eq:labor-success}. For the inequality, convexity of the exponential gives $\kappa e^0+(1-\kappa)e^{-y}\ge e^{-(1-\kappa)y}$ for $\kappa\in[0,1]$ and $y\ge0$; raise both sides to the power $n_K$.

(d) Expected value is $\varpi_pY_s+(1-\varpi_p)Y_f$ minus expected award payments. Raising $\varpi_p$ by $\Delta\varpi$ adds $\Delta\varpi(Y_s-Y_f)$. Each holder is still employed at cutover with probability $q_{\rm stay}=e^{-(\nu^q+(1-\kappa_{\rm ret})\nu^x)\tau^{\rm ann}}$, so the awards cost $n_KC^{\rm ret}q_{\rm stay}$ in expectation, which gives the condition. Write $y:=\nu^x\tau^{\rm ann}$ and $\mathfrak u(\kappa_{\rm ret}):=[\kappa_{\rm doc}+(1-\kappa_{\rm doc})e^{-(1-\kappa_{\rm ret})y}]^{n_K}$. Then $\mathfrak u(\kappa_{\rm ret})-\mathfrak u(0)=n_K(1-\kappa_{\rm doc})\kappa_{\rm ret}y+O(y^2)$ and $q_{\rm stay}=1+O(y)$. To first order in $y$ (the baseline exit $\nu^q\tau^{\rm ann}$ is of the same order), the condition reads $\varpi^0_pn_K(1-\kappa_{\rm doc})\kappa_{\rm ret}y(Y_s-Y_f)>n_KC^{\rm ret}$.

(e) Under a freeze, vacancies elsewhere arrive as a Poisson process with mean $\nu^qN^{\rm oth}T^w$ over the window. Each is filled by a displaced worker with probability $\varpi^{\rm rd}$, so the placements $N^{\rm pl}$ satisfy $\mathbb E N^{\rm pl}=\varpi^{\rm rd}\nu^qN^{\rm oth}T^w$. At most $N^{\rm dis}$ can be placed, and $\mathbb E\min\{N^{\rm dis},N^{\rm pl}\}\le\min\{N^{\rm dis},\mathbb EN^{\rm pl}\}$ by Jensen's inequality, since $\min\{N^{\rm dis},\cdot\}$ is concave. Dividing by $N^{\rm dis}$ gives the bound, which increases with $T^w$. $\square$

\paragraph{Proposition~\ref{prop:labor-talent}.}
(a) Only the labor share $l_{\rm AI}$ of the build is paid at $w^{\rm AI}$, so the build cost is $F_0(1-l_{\rm AI})+F_0l_{\rm AI}(1+\mathsf p)=F_0(1+l_{\rm AI}\mathsf p)$. The minimum horizon of Proposition~\ref{prop:tasks-transition}(b) is affine in $F$ with slope $1/(\varpi_p\bar b_pW)$.

(b) The capacity left after maintenance, $\Gamma-e_{\rm mnt}N$, builds at the rate $1/e_{\rm bld}$ processes per person-year. The linear equation $e_{\rm bld}\dot N=\Gamma-e_{\rm mnt}N$ with $N(0)=0$ has the solution $N(t)=\bar N_\Gamma(1-e^{-e_{\rm mnt}t/e_{\rm bld}})<\bar N_\Gamma$, which is increasing; inverting gives \eqref{eq:labor-TGamma}. With $y=N/\bar N_\Gamma$, $e_{\rm bld}N/\Gamma=ye_{\rm bld}/e_{\rm mnt}$, which gives the second form. The stretch factor is $-\ln(1-y)/y=\sum_{i\ge1}y^{i-1}/i$, a power series with positive coefficients, so it increases from 1 at $y=0^+$ and tends to infinity as $y\to1$.

(c) A hire starting at $t^{\rm s}\in[T^{\rm lead},T]$ has delivered, by $T$, $(T-t^{\rm s})-\delta^{\rm rp}\min\{T-t^{\rm s},T^{\rm rp}\}$ FTE-years. Integrating over starts at the rate $\dot\Gamma^{\rm hire}$ with $u=T-t^{\rm s}$ and $U:=T-T^{\rm lead}$ gives $\dot\Gamma^{\rm hire}\int_0^U[u-\delta^{\rm rp}\min\{u,T^{\rm rp}\}]\,du=\dot\Gamma^{\rm hire}[U^2/2-\delta^{\rm rp}T^{\rm rp}(U-T^{\rm rp}/2)]$ for $U\ge T^{\rm rp}$. Without the ramp, $\dot\Gamma^{\rm hire}U^2/2=e_{\rm bld}N$ gives $U=\sqrt{2e_{\rm bld}N/\dot\Gamma^{\rm hire}}$; the ramp lowers delivered capacity by $\dot\Gamma^{\rm hire}\delta^{\rm rp}T^{\rm rp}(U-T^{\rm rp}/2)$ and so delays completion.

(d) Hiring delivers the gap from $T^{\rm lead}$, at $w^{\rm AI}$ for $T^{\rm need}$ years, plus the ramp shortfall $\delta^{\rm rp}T^{\rm rp}$ (valued at the wage, since $T^{\rm need}\ge T^{\rm rp}$) and the hiring cost; the savings it gates wait $T^{\rm lead}$. Contracting pays $(1+u_{\rm ct})w^{\rm AI}$ for $T^{\rm need}$ and the savings wait $T^{\rm ct}$. Subtracting the two costs and dividing by $\Delta\Gamma w^{\rm AI}u_{\rm ct}>0$ gives the condition. As $T^{\rm need}\to\infty$ the condition fails, so hiring (or retraining) is cheaper for a permanent need. $\square$

\paragraph{Lemma~\ref{lem:labor-location}.}
By Proposition~\ref{prop:ext-timing}(d)(ii), an interior adoption date is $T^*=\ln[\pi_0/(c_0-c^Q-rF_0/v)]/\delta_\pi$. At the reference site the denominator is $c_0-c^Q_{\rm lab}-c^Q_{\rm loss}-rF_0/v=\bar c-c^{\rm fx}$; at site $\mathfrak c$ it is $\mathsf w_{\mathfrak c}c_0-\mathsf w_{\mathfrak c}c^Q_{\rm lab}-c^Q_{\rm loss}-rF_0/v=\mathsf w_{\mathfrak c}\bar c-c^{\rm fx}$, with the same $\pi_0$. Subtracting the two dates gives the lag. Moreover $\mathsf w_{\mathfrak c}(\bar c-c^{\rm fx})-(\mathsf w_{\mathfrak c}\bar c-c^{\rm fx})=c^{\rm fx}(1-\mathsf w_{\mathfrak c})\ge0$, so $(\bar c-c^{\rm fx})/(\mathsf w_{\mathfrak c}\bar c-c^{\rm fx})\ge1/\mathsf w_{\mathfrak c}$, with equality if and only if $c^{\rm fx}=0$ or $\mathsf w_{\mathfrak c}=1$. If $\mathsf w_{\mathfrak c}\bar c\le c^{\rm fx}$ the condition of Proposition~\ref{prop:ext-timing}(d)(ii) fails, and adoption never pays. $\square$

\paragraph{Proposition~\ref{prop:labor-residency}.}
(a) Readiness is the share of an activity's inputs that are machine-accessible (Section~\ref{sec:tasks-readiness}). Inputs unusable at $\mathfrak g$ are not accessible there, so readiness at $\mathfrak g$ is at most $(1-\omega^{\rm res}_j(\mathfrak g))\iota_j$. Feasibility then follows from the readiness floor of Proposition~\ref{prop:tasks-augment} and the capability test of Assumption~\ref{ass:ext-capability}, evaluated at the capability available in $\mathfrak g$, which is the frontier of $T^{\rm lag}_{\mathfrak g}$ years earlier. The permitted modes $\mathcal M_j$ of Definition~\ref{def:ext-effcost} contain an AI mode that uses the restricted inputs only if some region passes both tests, and $\mathcal J^{\rm per}$ in Definition~\ref{def:ext-universe} is defined by that set.

(b) With proportional restrictions, process readiness at $\mathfrak g$ is $(1-\omega)\iota_p$, where $\omega:=\omega^{\rm res}_p(\mathfrak g)$. By Proposition~\ref{prop:tasks-readiness}(b), full automation is optimal if and only if $(1-\omega)\iota_p\ge\iota^*_p$, that is, $\omega\le\bar\omega^{\rm res}_p$; otherwise no activity is automated and each runs in its best non-automated mode. By part (a) of that proposition, $W_p(\mathcal J_p;\cdot)$ is continuous and strictly increasing. If $\iota^*_p>\underline\iota_p$, continuity and the definition of $\iota^*_p$ give $W_p(\mathcal J_p;\iota^*_p)=0$, so $[W_p(\mathcal J_p;(1-\omega)\iota_p)]_+$ is continuous in $\omega$ and reaches zero at $\bar\omega^{\rm res}_p$. If $\iota^*_p=\underline\iota_p$, then $W_p(\mathcal J_p;\underline\iota_p)\ge0$, the value just below $\bar\omega^{\rm res}_p$ tends to it, and above $\bar\omega^{\rm res}_p$ full automation is infeasible, so the value drops by that amount. The last claim is part (c) of that proposition: full automation is feasible if and only if every activity's readiness is at least its floor.

(c) At region $\mathfrak g$ the unit cost of mode A consists of AI spend at the regional price, the quality cost at the readiness available there (Definition~\ref{def:ext-effcost}, with $\hat c^Q_j$ nonincreasing in readiness), the mechanism's annual cost and its expected replacement cost per unit, and costs common to all regions. A mechanism that lapses at the hazard $\nu^{\rm inv}$ and is replaced at cost $C^{\rm sw}$ each time costs $\nu^{\rm inv}C^{\rm sw}$ a year in expectation, the mean of a Poisson count of lapses. The best region minimizes the sum. The residency premium is nonnegative because the unrestricted case has full readiness, no premium and no mechanism costs. It raises the minimal AI-mode cost in \eqref{eq:ext-B} when mode A is the minimizer and leaves the non-AI modes unchanged. So it lowers $\mathcal B_j$ by at most $\Delta c^{\rm res}_j$, and by exactly that amount when no other AI mode is cheaper. $\square$

\paragraph{Corollary~\ref{cor:labor-pace}.}
If every wave lasts $T^{\rm eff}$, the program problem of Assumption~\ref{ass:ext-program} is unchanged except that $T_{\rm par}$ is replaced by $T^{\rm eff}$ in \eqref{eq:ext-progcost} and in the capacity limit per wave, so the proofs of Propositions~\ref{prop:ext-waves} and \ref{prop:ext-bandwidth} apply verbatim. The first two bounds are Proposition~\ref{prop:ext-bandwidth}(b) with $T^{\rm eff}$. For the third, if capacity satisfies $\Gamma(t)\le\Gamma_{\max}$, then $e_{\rm bld}\dot N\le\Gamma_{\max}-e_{\rm mnt}N$, and the comparison theorem for scalar differential equations gives $N(t)\le\bar N_{\Gamma_{\max}}(1-e^{-e_{\rm mnt}t/e_{\rm bld}})$. Hence reaching $N$ takes at least $T_{\Gamma_{\max}}(N)$, which is finite only if $N<\Gamma_{\max}/e_{\rm mnt}$. The survivor charge enters by Proposition~\ref{prop:labor-survivor}(b). $\square$
